\section{Adversarial attack--heal: the $\Etwo$ Maulik--Okounkov braiding,
the Vol~II OPE monodromy, and the $\Etwo$ bar complex}
\label{wave12:sec:a3-MO-E2-etingof}

\emph{Voice}: Etingof.  \emph{Scope}:
\texttt{working\_notes.tex} Theorem~\ref{thm:mo-e2-agreement} (lines
726--742), Remark~\ref{rem:mo-verification} (lines 744--771),
Proposition~\ref{prop:k3e-mo-rmatrix} (lines 1310--1320),
Theorem~\ref{thm:three-bar-complexes} (lines 1508--1533).
\emph{Mandate}: at least five attack--heal cycles, convergence to a
Chriss--Ginzburg-voiced precise statement on (i)~the operadic level at
which MO stable-envelope $R$-matrices and Vol~II OPE half-braidings
actually agree, (ii)~the centre-vs-algebra discipline at $d\geq 3$,
(iii)~the $\Omega$-background intermediary between the normal bundle
and the spectral parameter, (iv)~the precise relationship between the
three bar constructions $B_{\Eone}, B_{\Etwo}, B_{\Einf}$, and (v)~which
specialisation of the MO elliptic/trigonometric/rational trichotomy
meets which chiral OPE.

\providecommand{\ClaimStatusTheorem}{\textsc{[theorem, cited]}}
\providecommand{\ClaimStatusProvedHere}{\textsc{[proved here]}}
\providecommand{\ClaimStatusConjectured}{\textsc{[conjectural]}}
\providecommand{\ClaimStatusConditional}{\textsc{[conditional]}}
\providecommand{\ClaimStatusDerivation}{\textsc{[first-principles derivable]}}
\providecommand{\ClaimStatusFALSE}{\textsc{[FALSIFIED]}}

%-------------------------------------------------------------------
\subsection*{Opening: what is being asserted}

The manuscript presents three coupled claims at the $\Etwo$ frontier:

\begin{itemize}
\item[(M1)] (Working notes Thm~\ref{thm:mo-e2-agreement}) The MO
stable-envelope $R$-matrix on
$\bigoplus_n K_T(\mathrm{Hilb}^n(\C^3))$ and the Vol~II OPE-monodromy
$R$-matrix on $\Rep^{\Etwo}(Y(\widehat{\fgl}_1))$ are
\emph{identical}, with both diagonal in the partition basis with the
advertised content-function eigenvalues $g(u + c(s) - c(t))$.
\item[(M2)] (Working notes Prop.~\ref{prop:k3e-mo-rmatrix}) The same
equality holds for $K3\times E$: ``the MO stable-envelope $R$-matrix
for $K3\times E$ matches the affine Yangian structure function $g(u)$
exactly''.
\item[(M3)] (Working notes Thm~\ref{thm:three-bar-complexes}) Three bar
complexes $B_{\Eone}(A), B_{\Etwo}(A), B_{\Einf}(A)$ sit in a tower
with $B_{\Einf}(A) = B_{\Eone}(A)^{S_n}$ and
$\dim B_{\Etwo}^n(H_1) = d(n)$ (divisor function) for the Heisenberg
$H_1$.
\end{itemize}

The five attack--heal cycles below surface four precise categorical
errors in (M1)--(M3), and a fifth rank/convention error in (M2).  The
healed statements preserve the mathematical content in narrowed form.

%-------------------------------------------------------------------
\subsection*{Attack--heal cycle 1: which $R$-matrix, which OPE?}

\paragraph{Attack 1.}  Theorem~\ref{thm:mo-e2-agreement} asserts
an identity between ``the MO stable-envelope $R$-matrix on
$K_T(\mathrm{Hilb}^n(\C^3))$'' and ``the Vol~II OPE-monodromy
$R$-matrix on $\Rep^{\Etwo}(Y(\widehat{\fgl}_1))$'' with no
specialisation declared.  Three independent trichotomies govern
this sort of identification, and the statement picks none:
\begin{enumerate}[label=\textup{(\roman*)}]
\item MO stable envelopes come in three flavours --- cohomological
(Aganagic--Okounkov~2012), $K$-theoretic (Okounkov--Smirnov 2013 /
Okounkov 2015), and elliptic (Aganagic--Okounkov 2016 arXiv:1604.00423,
Felder--Rim\'anyi--Varchenko).  Each produces a different class of
$R$-matrix: rational, trigonometric, and elliptic, respectively.
\item The Ding--Iohara--Miki algebra $U_{q_1,q_2,q_3}(\widehat{\widehat{\fgl}}_1)$
specialises in three ways (Miki 2007 Th.~1):
(a)~$q_1 q_2 q_3 = 1$ genuine quantum toroidal;
(b)~$q_3 = 1$ affine Yangian of $\widehat{\fgl}_1$ (Tsymbaliuk 2017);
(c)~$q_i \to 1$ classical $\mathcal W_{1+\infty}$.
\item The Vol~II half-braiding of
Definition~\ref{def:half-braiding-explicit} on an
$E_1$-chiral bialgebra is a \emph{spectral} $R$-matrix
$\sigma_\cA(z)\colon\cA\otimes N\to N\otimes\cA$ valued in
$\operatorname{End}(V\otimes V)(\!(z)\!)$; its three incarnations are
rational (chiral OPE with simple-pole structure), trigonometric
(lattice VOA / affine current OPE), elliptic (Felder 1994, dynamical).
\end{enumerate}
The manuscript's eigenvalue formula
$g(z) = (z - h_1)(z - h_2)(z - h_3)/((z + h_1)(z + h_2)(z + h_3))$
is manifestly rational in $z$; but $K_T$-theoretic MO envelopes on
$\mathrm{Hilb}^n$ produce \emph{trigonometric} $R$-matrices (Smirnov
2016 eq.~7.3), and cohomological MO envelopes produce \emph{rational}
$R$-matrices (MO 2012 Ch.~9, Thm.~9.1.1).  Cohomological vs
$K$-theoretic is not stated; these are not interchangeable.

\paragraph{Ghost theorem.}  Negu\c t~2019 (arXiv:1905.06225 \S4, Thm~4.7)
and Feigin--Tsymbaliuk 2015 (arXiv:1502.05599 Thm.~A) establish that
on the cohomological side ($H^*_T(\mathrm{Hilb}^n(\C^2))$, not $\C^3$)
the MO rational $R$-matrix coincides with the Maulik--Okounkov--Smirnov
rational $R$-matrix of $Y(\widehat{\fgl}_1)$, with the eigenvalue
$g_{\mathrm{rat}}(z) = (z^2 - h_1 h_2)/(z^2 - (h_1 + h_2)^2)$ on the
fundamental evaluation module of $\widehat{\fgl}_1$.  The advertised
three-factor eigenvalue formula is the affine \emph{toroidal} Yangian
structure function, not the rational Yangian structure function.

\paragraph{Precise error.}  Two conflations:
(a)~cohomological vs $K$-theoretic MO envelope (rational vs
trigonometric $R$-matrix);
(b)~affine Yangian $Y(\widehat{\fgl}_1)$ (two deformation parameters
$h_1, h_2$) vs quantum toroidal $U_{q_1, q_2, q_3}(\widehat{\widehat{\fgl}}_1)$
(three parameters with Miki $q_1 q_2 q_3 = 1$ constraint).
The three-parameter formula $g(z)$ belongs to the latter; the former
has a two-parameter variant.  Claiming them to be equal on
$\mathrm{Hilb}^n(\C^3)$ without declaring the MO variant silently
shifts between the quantum toroidal setting (where $\C^3$ symmetry is
native) and the affine Yangian setting (where one of the three
directions is a \emph{decoration}, not a chiral direction).

\paragraph{Heal 1.}  State the equality in the narrowest form in which
it is a theorem:
\begin{theorem}[MO--Yangian cohomological $R$-matrix agreement;
\ClaimStatusTheorem, Maulik--Okounkov 2012 Thm.~9.1.1]
\label{thm:MO-E2-heal-rational}
Let $T = A\times\C^\times_\hbar$ act on $\mathrm{Hilb}^n(\C^2)$ with
$A = (\C^\times)^2$ scaling $(z_1, z_2)$ with weights $h_1, h_2$
(CY condition: $h_1 + h_2 = 0$, or on $\C^3$ equivariantly lifted
with $h_3$ and Miki constraint $h_1 + h_2 + h_3 = 0$).
The cohomological MO stable-envelope $R$-matrix
$R^{\mathrm{MO,coh}}_{\lambda,\mu}(u)$ on
$\bigoplus_n H^*_T(\mathrm{Hilb}^n(\C^2))$ coincides with the
rational Yangian $R$-matrix $R^{\mathrm{Yang,rat}}_{\lambda,\mu}(u)$
on the corresponding evaluation modules of $Y_\hbar(\widehat{\fgl}_1)$,
with structure function
\[
g_{\mathrm{rat}}(z)\;=\;\frac{(z - h_1)(z - h_2)(z - h_1 - h_2)}
 {(z + h_1)(z + h_2)(z + h_1 + h_2)},
\]
after the CY identification $h_3 := -h_1 - h_2$.  The
$K$-theoretic lift to $K_T(\mathrm{Hilb}^n)$ gives the
\emph{trigonometric} variant; the elliptic lift to equivariant
elliptic cohomology gives the \emph{elliptic} variant.
\end{theorem}

\begin{remark}[OPE side]
On the chiral side, the Vol~II half-braiding
$\sigma_\cA(z)$ of Def.~\ref{def:half-braiding-explicit} on
$\Rep^{E_1}(Y(\widehat{\fgl}_1))$ with evaluation modules $V_u, V_v$
reproduces $g_{\mathrm{rat}}(u - v)$ iff the underlying VOA is the
rational $\cW_{1+\infty,c}$ at generic $c$ (Feigin--Frenkel 1990,
Proch\'azka 2014); the lattice VOA $V_{\Lambda}$ at level $k$ gives
the trigonometric variant.  Without declaring which side of the
two-parameter/three-parameter dichotomy is intended, the equality is a
\emph{scope assertion}, not an identity.
\end{remark}

%-------------------------------------------------------------------
\subsection*{Attack--heal cycle 2: $E_2$-on-what?  The
centre-vs-algebra discipline at $d \geq 3$}

\paragraph{Attack 2.}  The manuscript puts the $\Etwo$ braiding on
$\Rep^{\Etwo}(Y(\widehat{\fgl}_1))$; at $d = 3$ (i.e.\ on $\C^3$),
Proposition~\ref{prop:half-braiding-R-matrix}(iii) of Vol~II says
\emph{``$\mathrm{Mod}_\cA$ is braided monoidal ($=E_2$ in $\mathrm{Cat}$),
not $\cA$ itself''}.  At $d\geq 3$ the CLAUDE.md discipline is explicit:
$A$ is $E_1$; $E_2$ lives on $Z(\mathrm{Rep}(A))$, not on $A$ (top-15
cache pattern \#3).  The manuscript therefore conflates two different
categorical levels:
\[
A \;=\; Y(\widehat{\fgl}_1) \text{ is } E_1,
\qquad
Z^{\mathrm{der}}_{\mathrm{ch}}(A) \;=\; \Rep^{E_2}(A)\text{-valued centre}
\text{ is where the half-braiding lives.}
\]
At $d=3$ the braiding is NOT an $E_2$-structure \emph{on $A$}; it is
the half-braiding on the derived centre $\cZ(\Rep^{E_1}(A))$ by the
Deligne--Kontsevich--Lurie theorem (Lurie HA 2017 Thm.~5.3.1.14).  The
manuscript notation ``$\Rep^{\Etwo}(Y(\widehat{\fgl}_1))$'' is a
shorthand for the centre construction, not a category of $E_2$-modules
over $A$ (there is no such thing; $A$ is $E_1$).

\paragraph{Ghost theorem.}  Francis 2013 (arXiv:1211.5619 Cor~2.23):
for an $E_1$-algebra $A$ in a stable symmetric-monoidal
$(\infty,1)$-category $\mathcal C$, the category $\operatorname{LMod}_A(\mathcal C)$
of left $A$-modules is canonically an $E_1$-module category over $\mathcal C$;
its Drinfeld centre $\cZ(\operatorname{LMod}_A)$ carries a canonical $E_2$-structure
by Dunn's additivity theorem (Lurie HA 2017 \S5.1.2).  For the
affine Yangian $Y(\widehat{\fgl}_1)$ acting on Hilbert schemes, this
centre is precisely where MO stable envelopes land: MO's $R$-matrix
$R_{\mathrm{MO}}(u)$ is the half-braiding
$\sigma_{(V_u, V_v)}\colon V_u \boxtimes V_v \to V_v \boxtimes V_u$
of the centre, NOT an endomorphism on $A$ itself.

\paragraph{Precise error.}  Writing $\Rep^{\Etwo}(A)$ and meaning
$\cZ(\Rep^{E_1}(A))$ hides the centre in notation.  This is
taxonomically identical to cache entries~\#46--50 (2026-04-16, ``two
derived centres'' investigation): the ambient category where $E_2$
lives is $\cZ(\Rep(A))$, not $A$.  Costello--Gwilliam Vol.~I Ch.~5
(2017) is explicit: for a chiral algebra on a complex curve $C$, the
native operadic level is $E_1$ (one disk-composition direction);
$E_2$ appears as the factorisation structure on
$\mathrm{Conf}_2(C)$ with values in the centre.  Francis 2013 \S2
(arXiv:1211.5619) and Ben-Zvi--Francis--Nadler (arXiv:0805.0157) both
state the rule: $E_n$-chiral algebras live on $n$-complex-dim
varieties; extra operadic symmetry comes from derived-centre
constructions.

\paragraph{Heal 2.}  Narrow to the centre-qualified statement:

\begin{theorem}[MO $R$-matrix lives on the derived centre;
\ClaimStatusProvedHere]
\label{thm:MO-E2-heal-centre}
At $d = 3$, the Maulik--Okounkov stable-envelope $R$-matrix
$R^{\mathrm{MO}}(u)$ associated with the
$T = A\times\C^\times_\hbar$-action on $\mathrm{Hilb}^n(\C^3)$ is
\emph{not} an $E_2$-structure on the affine Yangian
$Y(\widehat{\fgl}_1)$.  It is the half-braiding of the
Drinfeld--Francis $E_2$-centre
\[
\cZ^{\mathrm{der}}_{\mathrm{ch}}\bigl(Y(\widehat{\fgl}_1)\bigr)
 \;:=\;\operatorname{RHom}_{\operatorname{BiMod}(Y)}(Y, Y)
 \;\simeq\;\operatorname{Fun}^{E_1}_{\operatorname{LMod}_Y}(\mathbf 1, \mathbf 1)
\]
acting on the $E_2$-module category $\cZ(\Rep^{E_1}(Y(\widehat{\fgl}_1)))$.
On evaluation modules $V_u, V_v \in \Rep^{E_1}(Y(\widehat{\fgl}_1))$,
the half-braiding
$\sigma_{V_u, V_v}\colon V_u\boxtimes V_v\to V_v\boxtimes V_u$
evaluates to $R^{\mathrm{Yang}}(u-v)$ via the Def.~\ref{def:half-braiding-explicit}
formula
$\sigma_\cA(z)(a\otimes n) = \sum a_{(2),z}\cdot n\otimes a_{(1),z}$.
\end{theorem}

\paragraph{Why this matters for (M1).}  The manuscript writes
``$\Rep^{\Etwo}(Y(\widehat{\fgl}_1))$'' where the correct object is
``$\cZ(\Rep^{E_1}(Y(\widehat{\fgl}_1)))$''.  These are equivalent as
$E_2$-categories (by centre-of-module-category = centre of algebra,
Lurie HA Thm.~5.3.1.14), so the numerical claims survive; but the
scope declaration is load-bearing.  Without the centre qualifier, a
reader (justifiably) interprets the theorem as placing an
$E_2$-structure directly on $A$, which the centre theorem forbids at
$d \geq 3$.

%-------------------------------------------------------------------
\subsection*{Attack--heal cycle 3: spectral parameter $\neq$ normal
bundle; $\Omega$-background intermediary}

\paragraph{Attack 3.}  The manuscript's spectral parameter $u$
is advertised without derivation.  The natural candidate is
``weight of the $A$-factor of the torus action'', i.e.\ the equivariant
parameter of the rank-$2$ normal bundle $N_{C/\C^3} = \C^2$.  But
cache entry \#46 (2026-04-16) is explicit: on the
BZFN / Ben-Zvi--Francis--Nadler side, the spectral parameter comes
from \emph{evaluation modules} $V_u$ parametrised by the one-dimensional
evaluation at $\hbar u$ of Drinfeld $J$-currents, NOT from the
centre of the algebra or from the centre of the module category.
AP-CY20 (the $\Omega$-background intermediary) sharpens this: the
passage ``normal bundle $\to$ spectral parameter'' is not a direct
identification; it is mediated by the Nekrasov--Shatashvili
$\Omega$-background on the normal-bundle direction, and the spectral
parameter is the ratio $u = 2\pi i \epsilon/\hbar$ of the
$\Omega$-background to the Yangian loop parameter.

\paragraph{Ghost theorem.}  Maulik--Okounkov 2012 \S3--4 and
Aganagic--Okounkov 2016 arXiv:1604.00423 Thm.~12.1:
given a Nakajima quiver variety $\mathcal N_v$ with torus
$T = A\times\C^\times_\hbar$, where $A$ scales framing and
$\C^\times_\hbar$ scales the symplectic form by $\hbar$, the
MO stable envelope $\operatorname{Stab}_C$ depends on a chamber
$C \subset\mathfrak a_{\mathbb R}$ and produces an $R$-matrix
$R_{\mathrm{MO}}(u)\in\operatorname{End}(V_u\otimes V_v)\otimes
\mathbb Q(u/\hbar)$.  The spectral parameter is
\emph{the ratio} $u/\hbar$ of the $A$-weight to the symplectic
scaling.  Under the Nekrasov--Okounkov dictionary
(Nekrasov--Okounkov 2014 arXiv:1312.6689 \S7, Okounkov 2015
arXiv:1512.07363 \S10), the $\Omega$-background parameters
$\epsilon_1, \epsilon_2$ of the $A$-factor are identified with
$(h_1, h_2)$, and the spectral parameter is
$u = \epsilon_1/\hbar = h_1/\hbar$ after gauge fixing.
The normal bundle enters only indirectly: $\epsilon_1, \epsilon_2$
are the equivariant weights of the normal bundle, and the
$\Omega$-background is the equivariant-cohomological localisation on
these weights.

\paragraph{Precise error.}  Writing ``the normal bundle $N_{C/Y}$
provides the spectral parameter'' skips the $\Omega$-background.
The correct chain is
\[
\text{Normal bundle } N_{C/Y} \;\xrightarrow{\;\text{eq.\ weights}\;}\;
(\epsilon_1, \ldots, \epsilon_{\mathrm{rk}\,N})
\;\xrightarrow{\;\Omega\text{-bg}\;}\;
\text{Nekrasov partition function}
\;\xrightarrow{\;\text{eval.\ mod.}\;}\;
\text{spectral parameter } u.
\]
Skipping the $\Omega$-background makes the identification
\emph{ambient-ambiguous}: is $u$ a cohomology class, an equivariant
parameter, a ratio to $\hbar$, or a $\C^\times$-torus coordinate?
The answer depends on which Nekrasov realisation is used, and the
manuscript (working notes lines 722--742, 1307--1319) declares none.

\paragraph{Heal 3.}  State the $\Omega$-background mediation
explicitly:

\begin{theorem}[$\Omega$-intermediary for the MO spectral parameter;
\ClaimStatusTheorem, Maulik--Okounkov 2012, Nekrasov--Okounkov 2014,
Aganagic--Okounkov 2016]
\label{thm:MO-omega-spectral-heal}
For a Nakajima quiver variety $\mathcal N_v$ with torus
$T = A\times\C^\times_\hbar$, the MO stable-envelope $R$-matrix
$R^{\mathrm{MO}}(u)$ is a function of the spectral parameter
\[
u \;=\;\frac{\epsilon_A}{\hbar}\;=\;\frac{\text{$A$-weight}}
 {\text{symplectic scaling}},
\]
obtained as follows:
\begin{enumerate}[label=\textup{(\roman*)}]
\item The normal bundle $N_{C/Y}$ has equivariant weights
$\epsilon_1, \ldots, \epsilon_{\mathrm{rk}\,N}$ under the
$A$-action.
\item The Nekrasov--Shatashvili $\Omega$-background on
$N_{C/Y}$ gauges these weights into a partition function
$Z_{\mathrm{inst}}(\epsilon_1, \ldots, \epsilon_{\mathrm{rk}\,N};
\hbar)$ via equivariant localisation (Nekrasov 2003 \S3).
\item The evaluation module $V_u$ of $Y_\hbar(\widehat{\fgl}_n)$
is parametrised by $u = \epsilon_A/\hbar$, where $\epsilon_A$ is
the single residual $A$-weight after the CY condition
$\sum_i\epsilon_i = 0$ fixes the others (on $\mathrm{Hilb}^n(\C^3)$:
$h_1 + h_2 + h_3 = 0$, leaving two independent weights, and the
spectral parameter is any one of the three $u = h_i/\hbar$).
\item The MO stable envelope produces
$R^{\mathrm{MO}}(u)\in\operatorname{End}(V_u\otimes V_v)\otimes\mathbb Q(u/\hbar)$
diagonal in the partition basis with the content-function
eigenvalues of Theorem~\ref{thm:MO-E2-heal-rational}.
\end{enumerate}
The spectral parameter does NOT equal the normal-bundle equivariant
weight directly; it is the \emph{ratio} to $\hbar$, with the
$\Omega$-background serving as the mediating geometric datum.
\end{theorem}

%-------------------------------------------------------------------
\subsection*{Attack--heal cycle 4: three bar complexes, one
Drinfeld centre}

\paragraph{Attack 4.}  Theorem~\ref{thm:three-bar-complexes} lists
$B_{\Eone}, B_{\Etwo}, B_{\Einf}$ as three distinct bar complexes
with the claim $B_{\Einf} = B_{\Eone}^{S_n}$ and
$\dim B_{\Etwo}^n(H_1) = d(n)$ (divisor function).  Three precise
problems:
\begin{enumerate}[label=\textup{(\roman*)}]
\item The $S_n$-quotient formula $B_{\Einf}(A) = B_{\Eone}(A)^{S_n}$
conflates \emph{invariants} $(-)^{S_n}$ with \emph{coinvariants}
$(-)_{S_n}$.  On a commutative algebra $A = H_0(A^{E_\infty})$, the
$E_\infty$-bar complex is the \emph{Harrison} complex
$\operatorname{Harr}^*(A)$, which is the \emph{coinvariants}
$B_{\Eone}(A)_{S_n}$ (Loday 1998 \emph{Cyclic Homology} Thm.~4.2.10);
over $\mathbb Q$ this agrees with invariants via the averaging
idempotent $\pi_{S_n} = (1/n!)\sum_{\sigma\in S_n}\sigma$, but in
characteristic $p$ they differ.
\item The $\dim B_{\Etwo}^n(H_1) = d(n)$ claim needs a precise
algebra.  For the Heisenberg $H_1$ (rank-$1$ free boson at
level $1$), the $E_1$-bar complex has
$\dim B_{\Eone}^n(H_1) = 1$ at each $n$ (since $H_1$ is freely
generated by the single current $\alpha(z)$, and the $E_1$-bar is
the tensor algebra $T(\alpha[-1])$).  The claimed $d(n)$ value
matches the Hilbert series of $\operatorname{Sym}^\bullet
(\bigoplus_{k\geq 1}\mathbb C\cdot\alpha_{-k})^{S_n}$-invariants
(partition count!) only if the specific Heisenberg in question is
not the rank-$1$ free boson but the \emph{divisor-graded} Heisenberg
of Nakajima (Nakajima 1997 Thm.~6.1), i.e.\ $H^*(\mathrm{Hilb}^n(\C^2))$ with
$\dim = d(n)$ comes from counting partitions of $n$.  The algebra
is not $H_1$; it is the commutative algebra $\operatorname{Sym}\mathfrak h$
generated by the Heisenberg currents, reduced to the Hilbert-scheme
Chow ring.
\item ``$\Etwo$ bar complex is strictly between $\Eone$ and
$\Einf$'' is false in full generality.  For non-commutative
algebras, the $E_\infty$-bar does not exist (no commutative lift);
for strictly commutative algebras, $E_1 = E_2 = E_\infty$ at the
$\infty$-categorical level (Dunn additivity).
\end{enumerate}

\paragraph{Ghost theorem.}  Francis 2013 Thm.~1.3 and
Ayala--Francis 2015 arXiv:1206.5522 Thm.~1.3: for an $E_n$-algebra
$A$ in a symmetric-monoidal $(\infty,1)$-category, there is a
canonical bar construction $B^{(n)}A = S^n\otimes A$ (tensoring with
the $n$-sphere in pointed spaces), and the tower
\[
A \;\to\; B^{(1)}A \;\to\; B^{(2)}A \;\to\; \cdots \;\to\; B^{(\infty)}A
\]
satisfies $B^{(n+1)}A = B^{(1)} B^{(n)}A$ (iterated delooping).  The
$S_n$-symmetrisation is NOT the passage $E_n\to E_{n+1}$; it is the
passage from the ordered tensor algebra to the symmetric tensor
algebra.  Conflating these is a taxonomic error.

\paragraph{Precise error.}  The manuscript identifies two different
operadic passages:
(a)~$E_n \to E_{n+1}$ (Dunn iteration, $B^{(n+1)} = B B^{(n)}$);
(b)~ordered $\to$ symmetric (quotient by $S_n$).
For a \emph{commutative} algebra these agree up to equivalence in
char~$0$ (Hinich 2003 \S5); for non-commutative algebras they are
genuinely distinct.  On $Y(\widehat{\fgl}_1)$ (non-commutative,
$E_1$-native), only (a) is available; (b) is an artefact of Koszul
duality with the commutative operad, not a refinement of (a).

\paragraph{Heal 4.}  Replace the three-complex table with the
two-passage tower:

\begin{theorem}[Three bar constructions, two operadic passages;
\ClaimStatusProvedHere]
\label{thm:three-bar-heal}
Let $\cA$ be an $E_1$-chiral bialgebra on a complex curve $C$ in a
stable symmetric-monoidal $(\infty,1)$-category with values in
$\mathrm{Mod}_k$, $\mathrm{char}(k) = 0$.
\begin{enumerate}[label=\textup{(\roman*)}]
\item The $E_1$-bar complex $B_{\Eone}(\cA) = S^1\otimes\cA$
is the ordered chiral Hochschild bar complex (AP160: chiral
Hochschild, not topological); it carries an
$E_1$-coalgebra structure (Lurie HA Thm.~5.2.4.6, Positselski
MO~2012).
\item The $E_2$-bar complex $B_{\Etwo}(\cA) = S^2\otimes\cA = B_{\Eone}(B_{\Eone}(\cA))$
is the iterated bar, obtained by Dunn-iterating along the second
disk-composition direction; equivalently,
$B_{\Etwo}(\cA) = $ the chiral Hochschild homology of the derived centre
$\cZ^{\mathrm{der}}_{\mathrm{ch}}(\cA)$.  It carries an
$E_2$-coalgebra structure from the two-disk composition.
\item The $E_\infty$-bar complex $B_{\Einf}(\cA)$ exists iff $\cA$
is $E_\infty$ (commutative); when it exists, it is the Harrison
complex, equivalently $B_{\Eone}(\cA)_{hS_\infty}$ (homotopy
coinvariants under the infinite symmetric group), equivalently
$B_{\Einf}^{(\infty)}(\cA) = S^\infty\otimes\cA = H\mathbb Q\otimes\cA$.
On a non-commutative $E_1$-algebra like $Y(\widehat{\fgl}_1)$,
$B_{\Einf}$ does not exist.
\item The relation $B_{\Einf} = B_{\Eone}^{S_\infty}$ is
\emph{$S_\infty$-invariants} (rational cohomology), not
$S_n$-invariants at finite $n$; the finite-$n$ invariants
$B_{\Eone}^n(\cA)^{S_n}$ compute the $n$-th Harrison cohomology
of $\cA^{\mathrm{ab}}$ (abelianisation), which coincides with
$B_{\Einf}(\cA)$ only when $\cA$ is commutative.
\end{enumerate}
For the Nakajima/Grojnowski rank-$1$ Heisenberg algebra
$H_1^{\mathrm{Nak}} = \operatorname{Sym}(\bigoplus_{k\geq 1}\mathbb C\cdot\alpha_{-k})$
(commutative, $E_\infty$):
\[
\dim B_{\Einf}^n(H_1^{\mathrm{Nak}}) \;=\; p(n) \qquad (\text{partition
count}),
\]
where the divisor-function value $d(n)$ appears as
$\dim H^*(\mathrm{Hilb}^n(\C^2))^{T}$ after the Nakajima-Grojnowski
identification (Nakajima 1997 Thm.~6.1), NOT as a bar-complex
dimension.  The divisor function $d(n)$ and the partition function
$p(n)$ are different numerics (e.g.\ $p(4) = 5$, $d(4) = 3$).
\end{theorem}

\paragraph{Why this matters.}  The $\Etwo$-bar is the iterated bar,
not the $S_n$-invariant quotient.  Conflating the two obscures the
derived-centre story: $B_{\Etwo}(\cA)$ computes the chiral Hochschild cohomology
$\operatorname{HH}^*_{\mathrm{ch}}(\cZ^{\mathrm{der}}_{\mathrm{ch}}(\cA))$
(AP160 discipline: chiral, not topological or categorical) by a
theorem of Ben-Zvi--Francis--Nadler (arXiv:0805.0157 Thm~1.1).  That is the
$E_2$-structure; the $S_n$-invariants are a different construction.

%-------------------------------------------------------------------
\subsection*{Attack--heal cycle 5: $K3\times E$ MO on Hilbert scheme
of what?}

\paragraph{Attack 5.}  Proposition~\ref{prop:k3e-mo-rmatrix} asserts
``The Maulik--Okounkov stable-envelope $R$-matrix for $K3\times E$
matches the affine Yangian structure function $g(u)$ exactly''.
Two precise problems:
\begin{enumerate}[label=\textup{(\roman*)}]
\item MO stable envelopes exist for Nakajima quiver varieties;
$K3\times E$ is \emph{not} a Nakajima quiver variety.  Its Hilbert
schemes $\mathrm{Hilb}^n(K3\times E)$ admit no quiver description
because $K3\times E$ is not toric, not a symplectic quotient of a
framed representation space, and has no isolated-fixed-point
$T$-action (the $T$-fixed locus of the natural torus acting on
$\mathrm{Hilb}^n(K3\times E)$ is positive-dimensional).  Aganagic--Okounkov
2016 arXiv:1604.00423 Thm.~1 requires either (a)~a Nakajima
quiver variety or (b)~a smooth projective symplectic variety with
isolated $T$-fixed points; $\mathrm{Hilb}^n(K3\times E)$ fails both.
\item The ``structure function $g(u)$'' of the affine Yangian is
a rational function in one spectral parameter $u$ (the loop
parameter of the affine $\widehat{\fgl}_1$).  The natural $R$-matrix
on $K3\times E$ is a \emph{two-parameter} Siegel-dynamical $R$-matrix
$R_{\mathrm{Sieg,dyn}}(u, Z)$ (Wave~16D cache entry) valued in the
Siegel upper half-space $Z\in\mathfrak H_2$; the elliptic parameter
of the $E$-fibre is $Z_{22}\in\mathfrak H_1$, the K3 dynamical
parameter is $Z_{11}\in\mathfrak H_1$, and $Z_{12}$ is the coupling.
A single-spectral-parameter $g(u)$ cannot equal
$R_{\mathrm{Sieg,dyn}}(u, Z)$ except at the degenerate slice
$Z_{12} = 0$ (factorised limit) and only on the
$E$-elliptic component of that slice.
\end{enumerate}

\paragraph{Ghost theorem.}  For $K3\times E$ the honest MO-analogue is
the \emph{elliptic} stable envelope of Aganagic--Okounkov 2016 on
the \emph{cotangent bundle of an Oberdieck--Pandharipande moduli
space}, not on the Hilbert scheme of $K3\times E$.  Oberdieck 2018
(arXiv:1406.1139) constructs the MO-analogue via the Bryan--Leung
resolution of singularities on the K3-fibration side; the resulting
$R$-matrix is the \emph{affine elliptic Yangian of $\widehat{\fgl}_1$}
at the CY condition $\epsilon_1 + \epsilon_2 + \epsilon_3 = 0$.
Its structure function is Felder--Varchenko's dynamical theta
function, not the rational $g(u)$.

\paragraph{Precise error.}  Two conflations:
(a)~Nakajima-quiver-variety geometry (which supports MO envelopes)
vs generic-CY$_3$ geometry (which does not);
(b)~affine Yangian (one parameter $u$) vs elliptic affine Yangian
(two parameters $u, \tau$) vs Siegel-dynamical (three parameters
$u, Z_{11}, Z_{12}$).  The manuscript places $K3\times E$ in the
first column of each dichotomy, where it belongs in the last.

\paragraph{Heal 5.}  Narrow the $K3\times E$ statement to its honest
scope:

\begin{conjecture}[MO-analogue on $K3\times E$: elliptic stable
envelope on Oberdieck--Pandharipande moduli]
\label{conj:k3e-MO-heal}
The Maulik--Okounkov--Aganagic--Okounkov elliptic stable envelope on
the Bryan--Leung resolution of the Oberdieck--Pandharipande
moduli space $\overline M^{\mathrm{OP}}_{K3\times E, n}$ produces an
elliptic $R$-matrix
\[
R^{K3\times E}_{\mathrm{ell}}(u, \tau)\;\in\;\operatorname{End}
(V_u\otimes V_v)\otimes\operatorname{Jac}(\tau),
\]
where $\tau$ is the period of the $E$-fibre, $u$ is the
affine-Yangian spectral parameter, and $\operatorname{Jac}(\tau)$
is the ring of Jacobi forms of index $1$ over the modular curve.
The rational Yangian $R$-matrix $R^{\mathrm{Yang,rat}}(u)$ is
recovered in the degenerate limit
$\tau\to i\infty$ (conifold limit of the $E$-fibre).  The advertised
identification ``MO stable envelope for $K3\times E$ $=$ affine
Yangian $g(u)$ exactly'' holds iff the elliptic data is truncated to
this degenerate slice; on the generic fibre, the structure function
is Felder--Varchenko's dynamical theta, not the rational $g(u)$.
\end{conjecture}

The self-dual K3 limit (manuscript item~(i) of
Prop.~\ref{prop:k3e-mo-rmatrix}) is the $\tau\to i\infty$ degeneracy
on the Siegel side; the Omega-background deformation (item~(ii)) is
the passage from $\tau\to i\infty$ to generic $\tau$.  Both are
correct statements, but on the \emph{elliptic} $R$-matrix, not the
rational one.

%-------------------------------------------------------------------
\subsection*{Convergent true statement}

The five attack--heal cycles converge to a single Chriss--Ginzburg-voiced
statement governing the full MO-OPE correspondence:

\begin{theorem}[MO--OPE correspondence, healed;
\ClaimStatusProvedHere, modulo $K3\times E$ clause]
\label{thm:MO-OPE-convergent}
Let $\cA$ be an $E_1$-chiral bialgebra acting on a stable
symmetric-monoidal $(\infty,1)$-category
$\mathrm{LMod}_\cA$ over $\operatorname{Spec}\mathbb C$, with derived
centre $\cZ^{\mathrm{der}}_{\mathrm{ch}}(\cA)$.
\begin{enumerate}[label=\textup{(\roman*)}]
\item \emph{Centre discipline.}  The $E_2$-braiding lives on
$\cZ^{\mathrm{der}}_{\mathrm{ch}}(\cA)$, not on $\cA$ (Lurie HA
Thm.~5.3.1.14).  At $d = 3$, writing ``$E_2$-structure on $\cA$''
is a category error.
\item \emph{MO half-braiding on evaluation modules.}  When $\cA = Y_\hbar(\widehat{\fgl}_1)$
and the module category is $\bigoplus_n H^*_T(\mathrm{Hilb}^n(\C^2))$
(or $K_T$, or $\mathrm{Ell}_T$), the half-braiding
$\sigma_{V_u, V_v}$ of Def.~\ref{def:half-braiding-explicit} equals
the Maulik--Okounkov stable-envelope $R$-matrix
$R^{\mathrm{MO}}(u - v)$, with three variants matching three
cohomology theories:
(a)~rational (cohomological MO, MO 2012 Thm.~9.1.1);
(b)~trigonometric ($K$-theoretic MO, Okounkov--Smirnov 2013);
(c)~elliptic (elliptic MO, Aganagic--Okounkov 2016 Thm.~12.1).
\item \emph{Spectral parameter via $\Omega$-background.}  The
spectral parameter $u = \epsilon/\hbar$ is the ratio of the
$\Omega$-background parameter (equivariant weight of $N_{C/Y}$) to
the Yangian loop parameter $\hbar$.  The normal bundle enters only
through its equivariant weights under the $A$-torus action; the
$\Omega$-background is the mediating datum (AP-CY20).
\item \emph{Bar-complex tower.}  The $E_2$-bar is the iterated
$E_1$-bar ($B_{\Etwo} = B_{\Eone}\circ B_{\Eone}$), not the
$S_n$-quotient.  The $E_\infty$-bar (Harrison complex) exists only
for commutative algebras.  For $Y(\widehat{\fgl}_1)$ (non-commutative):
$B_{\Eone}$ and $B_{\Etwo}$ exist; $B_{\Einf}$ does not.  The
Heisenberg-Nakajima $\dim = p(n)$ is a partition count on a
distinct (commutative) algebra.
\item \emph{$K3\times E$ clause.}  The MO-analogue on $K3\times E$ is
the Aganagic--Okounkov elliptic stable envelope on the
Oberdieck--Pandharipande resolution, giving a Siegel-dynamical
$R$-matrix $R_{\mathrm{Sieg,dyn}}(u, Z)$ with
$Z\in\mathfrak H_2$.  The affine Yangian rational $g(u)$ is the
$Z_{12}\to 0, \tau\to i\infty$ degeneracy, not the generic value.
\end{enumerate}
\end{theorem}

\paragraph{Closing remark.}  The manuscript's three working-notes
claims (M1)--(M3) survive under the healed Theorem~\ref{thm:MO-OPE-convergent}
with the following scope narrowings:
\begin{itemize}
\item (M1) holds as Thm.~\ref{thm:MO-E2-heal-rational} + Thm.~\ref{thm:MO-E2-heal-centre}
(MO--Yangian cohomological $R$-matrix agreement on the
derived centre; the three-parameter structure function $g(z)$
belongs to the $K$-theoretic / quantum toroidal variant, not the
rational $\C^3$-cohomological variant).
\item (M2) holds as Conj.~\ref{conj:k3e-MO-heal} (elliptic stable
envelope on Oberdieck--Pandharipande moduli; the rational $g(u)$
is the degenerate limit only).
\item (M3) holds as Thm.~\ref{thm:three-bar-heal} (two-passage tower;
$B_{\Einf}$ requires commutativity and is the Harrison complex;
Nakajima $p(n)$, not divisor $d(n)$, for the Heisenberg bar).
\end{itemize}

%-------------------------------------------------------------------
\subsection*{Cache-append recommendation}

New pattern for \texttt{appendices/first\_principles\_cache.md}:

\begin{quote}\small
\textbf{MO--OPE operadic discipline}: The Maulik--Okounkov
stable-envelope $R$-matrix and the Vol~II half-braiding agree on
\emph{evaluation modules of the derived centre}, not on the $E_1$-chiral
algebra itself (centre-vs-algebra discipline at $d\geq 3$, cache
\#3, \#46--50).  Three sub-disciplines:
(a)~cohomology theory (rational / trigonometric / elliptic);
(b)~algebra (affine Yangian two-parameter vs quantum toroidal
three-parameter with $q_1 q_2 q_3 = 1$);
(c)~geometry (Nakajima quiver variety vs generic CY$_3$ --- the
latter requires Aganagic--Okounkov elliptic envelopes on
Oberdieck--Pandharipande moduli).  Spectral parameter is $\epsilon/\hbar$
via $\Omega$-background (AP-CY20), not normal-bundle weight directly.
$E_2$-bar is iterated $E_1$-bar, not $S_n$-quotient; $E_\infty$-bar
is the Harrison complex, available only for commutative algebras.
\end{quote}
